During the early stages of recursive synthetic training, spectral indicators of the embedding space suggest the model is improving, while downstream performance degrades.
We characterize this disconnect by tracking spectral properties across 10 recursive generations of Pythia-410M.
The result is a reproducible three-phase pattern consistent across three seeds: spectral decompression (generations 0--2), where spectral log-volume (SLV) increases while HellaSwag drops 3.3 percentage points; collapse onset (generations 2--4), where spectral indicators reverse; and accelerating collapse (generations 4--10).
Of eight candidate indicators, SLV is the most effective indicator: it requires one forward pass (${<}\,1$\,s), captures changes across all spectral directions, and declines monotonically after the decompression peak.
Varying the mixing ratio across five levels yields a trajectory map localizing the critical mixing ratio $r_{\text{crit}} \in (0.25, 0.3]$, separating monotonic expansion from a transitional regime with initial contraction then recovery.
A proportional controller reveals a prevention--detectability tradeoff: real-data mixing suppresses spectral contraction, so the rate-based signal vanishes precisely when prevention succeeds---a particularly strong form of the ecological intervention--warning tradeoff.
Setting $r > r_{\text{crit}}$ prevents collapse without per-generation monitoring.
The three-phase trajectory reproduces on GPT-2-medium (355M), providing preliminary cross-architecture evidence\footnote{Code and data are available at \url{https://anonymous.4open.science/r/spectral_collapse_ews-1B2C/}.}.
